% ---- pg-187 ----
表中的 $X_{1},X_{2},X_{3}$ 分别代表社会科学与历史、语言表达及自然科学三门课程的考试成绩. 由表中的数据可得
\[
\bar{\mathbf{x}}=\begin{bmatrix}526.29\\54.69\\25.13\end{bmatrix}, \quad \mathbf{S}=\begin{bmatrix}5\,808.06&597.84&222.03\\597.84&126.05&23.39\\222.03&23.39&23.11\end{bmatrix}
\]

现计算 $\mu_{1},\mu_{2},\mu_{3}$ 的 95\% 联合置信区间. 由于
\[
\frac{p(n-1)}{n-p}F_{p,n-p}(\alpha)=\frac{3\times(87-1)}{(87-3)}F_{3,84}(0.05)=\frac{3\times86}{84}\times2.7=8.29
\]

可以得到联合置信区间[见\hyperref[eq:5-24]{式(5-24)}]为
\[
526.29-\sqrt{8.29}\times\sqrt{\frac{5\,808.06}{87}}\le \mu_{1}\le 526.29+\sqrt{8.29}\times\sqrt{\frac{5\,808.06}{87}}
\]

即
\[
503.06\le \mu_{1}\le 550.12
\]

\[
54.69-\sqrt{8.29}\times\sqrt{\frac{126.05}{87}}\le \mu_{2}\le 54.69+\sqrt{8.29}\times\sqrt{\frac{126.05}{87}}
\]

即
\[
51.22\le \mu_{2}\le 58.16
\]

\[
25.13-\sqrt{8.29}\times\sqrt{\frac{23.11}{87}}\le \mu_{3}\le 25.13+\sqrt{8.29}\times\sqrt{\frac{23.11}{87}}
\]

即
\[
23.65\le \mu_{3}\le 26.61
\]

除语言表达成绩外, 其边缘 $Q$-$Q$ 图及二维散点图表明, 并无迹象表明有严重背离大学资格考试数据的正态性的现象产生(见练习 5.18). 此外, 即使数据不是严格服从正态分布, 这么大的样本容量也足以保证该方法的合理性(见 5.5 节).

由于上述三个置信区间均必须以 95\% 的置信度成立, 因此其联合 $T^{2}$ 区间要比一元置信区间更宽些. 又因为可以用同样的置信度来构造它们之间差值的置信区间, 所以这些置信区间也许比通常所需的更宽.

例如, 对 $\mathbf{a}'=[0,1,-1]$, $\mu_{2}-\mu_{3}$ 的置信区间的两个端点为
\[
(\bar{x}_{2}-\bar{x}_{3})\pm\sqrt{\frac{p(n-1)}{(n-p)}F_{p,n-p}(0.05)}\sqrt{\frac{s_{22}+s_{33}-2s_{23}}{n}}
\]
\[
=(54.69-25.13)\pm\sqrt{8.29}\times\sqrt{\frac{126.05+23.11-2\times23.39}{87}}
\]
\[
=29.56\pm3.12
\]

故 $\mu_{2}-\mu_{3}$ 的 95\% 置信区间为 $(26.44,32.68)$. 对其他差值同样可构造其联合置信区间.

最后, 能以相同的 95\% 置信度构造每对均值的置信椭圆. 比如, 对 $(\mu_{2},\mu_{3})$ 这对均值来说, 有
\[
87\times[54.69-\mu_{2}, 25.13-\mu_{3}]\begin{bmatrix}126.05&23.39\\23.39&23.11\end{bmatrix}^{-1}\begin{bmatrix}54.69-\mu_{2}\\25.13-\mu_{3}\end{bmatrix}
\]
\[
=0.849\times(54.69-\mu_{2})^{2}+4.633\times(25.13-\mu_{3})^{2}
\]
\[
\qquad -2\times0.859\times(54.69-\mu_{2})\times(25.13-\mu_{3})\le 8.29
\]

该椭圆及另外两对均值的 95\% 置信椭圆的图形见\hyperref[fig:5.3]{图 5.3}. 图中还给出了这些椭圆在其均值分量轴上的投影, 这些投影即为 $T^{2}$ 区间.
